\section{蕴含与充要条件}
\label{sec:01-Mathematical-Language/01-Logic/02}

你朋友说：``要是期末每科都上 90，总评就是 A。''

期末你拿了 A。但有一科是 89，老师给四舍五入了。你朋友那个条件到底算不算被推翻？直觉告诉你：他没说拿 A 就一定每科 90 以上，所以你拿了 A 也反驳不了他。反过来，如果你每科都 92 却只拿了 B，那他倒确实说错了。

这个场景里的不对称，就是本节要拆开来看的东西。上一节我们用真值表摆平了 \(P\to Q\) 的四种真假组合，但真值表只管``什么时候算对、什么时候算错''，不管这条箭头在推理中到底怎么用。同一根箭头，从左边往右看和从右边往左看，说的是两件完全不同的事——而你每天读定理、写证明的时候，都得在这两个视角之间来回切换。

\subsection{一条箭头，两个读法}

先固定一个具体的蕴含：

\[
n\text{ 能被 4 整除}\quad\longrightarrow\quad n\text{ 是偶数}.
\]

拿 \(n=12\) 试一下：12 是 4 的倍数，12 是偶数，没出问题。拿 \(n=6\)：6 不是 4 的倍数，但 6 是偶数。这条蕴含有没有说 6 不对？没有——它只关心``4 的倍数''那批数字是不是偶数，不碰根本不满足前提的数。拿 \(n=3\)：既不是 4 的倍数也不是偶数，一样不构成反例。

唯一能推翻它的，是一个能被 4 整除的奇数。找一圈，没有。所以蕴含成立。

现在换个角度看同一根箭头。你有一个数字，想知道它是不是偶数。如果人家告诉你``它是 4 的倍数''，那你可以立刻下结论：是偶数。这个方向叫\textbf{充分条件}——有了 \(P\) 就足够了。但反过来，人家只告诉你``它是偶数''，你能不能说它一定是 4 的倍数？\(n=2\) 是偶数但不是 4 的倍数，不行。所以``是偶数''只是结论成立时\textbf{必须}满足的条件——不满足偶数的话，4 的倍数根本没戏——但它本身不足以让你下结论。这个方向叫\textbf{必要条件}。

记住这一张图就够了：

\[
P\quad\longrightarrow\quad Q.
\]

箭头的尾巴是充分的，箭头的尖是必要的。同一条箭头，换个读法而已，不是两条独立的规律。

微积分里最经典的例子：可微一定连续，连续不一定可微。\(f(x)=|x|\) 在 \(x=0\) 连续，但从左边趋近的斜率是 \(-1\)，右边是 \(+1\)，尖角卡在那里，导数不存在。所以``可微''是``连续''的充分条件，``连续''是``可微''的必要条件。知道一个函数可微，你可以放心说它连续；只知道它连续，不能反推可微。

\subsection{``若''、``仅当''与箭头方向}

自然语言坑就坑在这里。下面四句话，甩给没练过的读者，至少会卡两下：

\begin{itemize}
\tightlist
\item
  ``若 \(P\)，则 \(Q\)'' \(\longrightarrow\) \(P\to Q\)（这个不意外）；
\item
  ``\(P\) 仅当 \(Q\)'' \(\longrightarrow\) 还是 \(P\to Q\)（第一次见多半愣住）；
\item
  ``只要 \(P\)，就 \(Q\)'' \(\longrightarrow\) 依然是 \(P\to Q\)；
\item
  ``\(P\) 当且仅当 \(Q\)'' \(\longrightarrow\) 双向都成立。
\end{itemize}

``仅当''为什么箭头是 \(P\to Q\) 而不是 \(Q\to P\)？你试着把句子拆开。``\(P\) 仅当 \(Q\)''的意思是：\(P\) 成立这件事，只有在 \(Q\) 也成立的时候才可能发生。换句话说，\(Q\) 不成立的话，\(P\) 一定不成立。那就是 \(\neg Q\to\neg P\)，逆否回去就是 \(P\to Q\)。

碰到拿不准的中文句子，不要靠语感猜。直接问自己：\textbf{什么情况能推翻这句话？}比如``只有注册用户才能评论''——它禁止的是什么？一个没注册却评论了的人。所以它的逻辑形式是：能评论 \(\to\) 已注册。箭头从``能评论''指向``已注册''。能评论是已注册的充分条件？不对——能评论意味着已注册，所以``已注册''是必要条件。每句话都把这个测试跑一遍，方向慢慢就内化了。

\subsection{逆命题、否命题与逆否命题}

给定 \(P\to Q\)，可以机械地扭出三个变体：

\begin{itemize}
\tightlist
\item
  逆命题：\(Q\to P\)（箭头掉头）；
\item
  否命题：\(\neg P\to\neg Q\)（两边各自否定，箭头方向不动）；
\item
  逆否命题：\(\neg Q\to\neg P\)（箭头掉头且两边否定）。
\end{itemize}

还是用``能被 4 整除 \(\to\) 是偶数''来试：

\begin{itemize}
\tightlist
\item
  逆命题``偶数都能被 4 整除''——\(n=2\) 就倒了。
\item
  否命题``不能被 4 整除就不是偶数''——\(n=6\) 不能被 4 整除，但它是偶数，又倒了。
\item
  逆否命题``不是偶数的数不能被 4 整除''——奇数确实不可能是 4 的倍数，成立。
\end{itemize}

你大概注意到了：原命题成立，逆否命题也跟着成立；否命题和逆命题也绑在一起。这不是巧合。列个真值表就能验出来，\(P\to Q\) 和 \(\neg Q\to\neg P\) 在所有四种真假组合下结果一模一样；而否命题 \(\neg P\to\neg Q\) 其实就是逆命题的逆否命题，所以它俩自动等价。

这个事实有一个直接的用处：\textbf{逆否证法}。有时候正面从 \(P\) 推到 \(Q\) 不好下手，但从 \(\neg Q\) 推到 \(\neg P\) 反而顺畅。你做的事是在证明原命题，只是换了个姿势。

什么时候该走逆否？一条简单的判断标准：看 \(\neg Q\) 是不是比 \(P\) 更容易处理。比如你要证``若 \(n^2\) 是偶数，则 \(n\) 是偶数''。正面来？设 \(n^2=2k\)，想从中解出 \(n\) 的奇偶性，不顺手。走逆否：假设 \(n\) 是奇数，\(n=2k+1\)，那 \(n^2=4k^2+4k+1=2(2k^2+2k)+1\)，奇数。逆否成立，原命题证完。``\(\neg Q\) 好不好处理''就是开关。

\subsection{充要条件需要两个方向}

\(P\leftrightarrow Q\) 没有什么神秘的新机制，它就是两条蕴含拧在一起：

\[
(P\to Q)\land(Q\to P).
\]

证``当且仅当''的时候，两个方向都得老老实实走完，一个都不能省。写完一个方向就收工，跟只做了半道题没有区别。

看一个标准的充要条件：对整数 \(n\)，

\[
n\text{ 为偶数}\quad\Longleftrightarrow\quad n^2\text{ 为偶数}.
\]

正向（\(\Rightarrow\)）：设 \(n=2k\)，那 \(n^2=4k^2=2(2k^2)\)，偶数。干净。

反向（\(\Leftarrow\)）：前面刚用逆否证法走过一遍，再来一次。假设 \(n\) 是奇数，\(n=2k+1\)，则 \(n^2=2(2k^2+2k)+1\)，奇数。所以 \(n^2\) 为偶数时 \(n\) 不可能是奇数，只能是偶数。反向证完。

两个方向的证明手法可以完全不同：正向是直接构造，反向用了逆否。关键是你得让读者一眼看出哪一段负责哪条箭头。写的时候在段落开头扔一句``（\(\Rightarrow\)）''或``（\(\Leftarrow\)）''就足够清楚了。

\subsection{定理条件不是装饰品}

读到一个定理的时候，别只盯着结论。把论域、假设和结论三个部件拆开，然后追问一句：每条假设在防什么？

拿最值定理练手：

\begin{theorem}[最值定理]
若 \(f\) 在闭区间 \([a,b]\) 上连续，则 \(f\) 在该区间上取得最大值和最小值。
\end{theorem}

假设有两个：``闭区间''和``连续''。缺了任何一条，结论就塌了。

把``闭区间''换掉。\(f(x)=x\) 在开区间 \((0,1)\) 上连续，但它的值可以无限接近 1，却永远碰不到 1——没有最大值。同样，你也没法说最小值是 0，因为 0 不在区间里。闭区间的作用是把两个端点关在里面，不让函数值溜到边界外却取不到。

把``连续''换掉。在 \([0,1]\) 上定义一个函数：\(f(x)=x\) 对 \(0\le x<1\)，而 \(f(1)=0\)。它不连续，在 \(x=1\) 处跳崖了。函数值在区间上可以任意接近 1（比如 \(f(0.999)=0.999\)），但没有一个点真的取到 1——因为接近 1 的那些 \(x\)，函数值是它们自己，而唯一可能取 1 的 \(x=1\)，函数值却是 0。最大值的候选被不连续处吞掉了。

定理的条件是``结论不翻车所需的最低结构''。但注意：证明里用到某个条件，不等于这个条件是结论的\textbf{必要}条件。有可能存在更厉害的定理，用更弱的假设也能推出同样的结论。读书的时候要区分``论证在哪个步骤调用了条件''和``结论在逻辑上是不是离了条件就活不了''，这是两笔账。

\begin{remark}
很多教科书只说``闭区间连续函数必有最值''，不给反例。但反例才是真正帮你记住条件的东西。下次看到一个定理，给自己 30 秒：随便删掉一条假设，动手捏一个反例出来。捏不出来？说明你还没理解那条假设在干什么。
\end{remark}

\subsection{反例的枪口要对准 \(P\land\neg Q\)}

推翻 \(P\to Q\)，只需要一个同时满足 \(P\) 和 \(\neg Q\) 的对象。注意：必须同时满足。随便甩一个跟命题无关的例子是打不中的。

比如命题``所有素数都是奇数''。你说 9 不是素数但它是奇数——这根本不算反例。9 不满足前提``是素数''，连上场资格都没有。真正的反例是 2：它是素数（满足 \(P\)），但不是奇数（满足 \(\neg Q\)）。一枪命中。

另一个常见错误：拿一个既不满足 \(P\) 也不满足 \(Q\) 的东西出来。比如反驳``能被 4 整除则是偶数''，你举 3。3 不能被 4 整除，也不是偶数——这能说明什么？什么都不说明。反例不是随便找个不顺眼的数字，它必须精准地坐在 \(P\) 的真值区域里，却在 \(Q\) 的区域外。

检验反例的三步：
\begin{itemize}
\tightlist
\item
  它满足前提 \(P\) 吗？
\item
  它不满足结论 \(Q\) 吗？
\item
  以上两条同时成立吗？
\end{itemize}

三步都通过，才是一个合格的杀手。

\subsection{箭头没有告诉你的东西}

蕴含告诉你的是：只要 \(P\) 真，\(Q\) 就不可能假。但它不告诉你以下任何一件事：

\begin{itemize}
\tightlist
\item
  \(P\) 是 \(Q\) 的\textbf{原因}。逻辑蕴含和因果关系是两套系统。``下雨则地湿''的箭头成立，但箭头本身不说雨是地湿的原因——它只说找不到下雨而地不湿的情况。
\item
  如果 \(Q\) 真，\(P\) 怎么样。从 \(P\to Q\) 和 \(Q\) 真，推不出 \(P\) 真。地湿不一定是因为下雨。
\item
  如果 \(P\) 假，\(Q\) 怎么样。从 \(P\to Q\) 和 \(\neg P\)，也推不出 \(\neg Q\)。没下雨，地也可能因为洒水车而湿。
\item
  在真实世界里，``如果 \(P\) 则 \(Q\)''是否仍然成立。逻辑只管真值指派，不管世界怎么变。今天``下雨则地湿''成立，明天修了个棚子，可能就不成立了。
\end{itemize}

这些都是推理中容易踩的坑：肯定后件、否定前件、把相关性当蕴含。每踩一次，回来翻翻这一节，坑会慢慢变浅。
